\clearpage
\oldpage[II]{207}
\markboth{제II장}{제II장}
\phantomsection
\addcontentsline{toc}{section}{기호 색인}
\section*{기호 색인}
\thispagestyle{plain}

\begingroup
\small
\setlength{\parindent}{0pt}
\setlength{\parskip}{2pt}
\newcommand{\egaShProj}{\mathcal{P}\!\operatorname{roj}}

$\mathcal{A}(X),\mathcal{A}(\mathcal{F}),\mathcal{A}(\mathcal{B})$
($X$는 $S$-준스킴, $\mathcal{F}$는 $\mathcal{O}_X$-가군,
$\mathcal{B}$는 $\mathcal{O}_X$-대수): \hyperref[II.1.1.1-ko]{1.1.1}.

$\mathcal{A}(h)$ ($h$는 $S$-사상): \hyperref[II.1.1.2-ko]{1.1.2}.

$\Spec(\mathcal{B})$ ($\mathcal{B}$는 $\mathcal{O}_S$-대수): \hyperref[II.1.3.1-ko]{1.3.1}.

$\widetilde{\mathcal{M}}$ ($\mathcal{M}$은 $\mathcal{B}$-가군,
$\mathcal{B}$는 $\mathcal{O}_S$-대수): \hyperref[II.1.4.3-ko]{1.4.3}.

$\mathbf{T}(E),\mathbf{S}(E),\mathbf{S}_A(E)$ ($E$는 $A$-가군): \hyperref[II.1.7.1-ko]{1.7.1}.

$\mathbf{S}(\mathcal{E}),\mathbf{S}_{\mathcal{A}}(\mathcal{E})$
($\mathcal{E}$는 $\mathcal{A}$-가군): \hyperref[II.1.7.4-ko]{1.7.4}.

$\mathbf{V}(\mathcal{E})$ ($\mathcal{E}$는 준연접 $\mathcal{O}_S$-가군): \hyperref[II.1.7.8-ko]{1.7.8}.

$S_n,S_+,M_n,S^{(d)},M^{(d,k)},M^{(d)}$
($S$는 등급환, $M$은 등급 $S$-가군): \hyperref[II.2.1.1-ko]{2.1.1}.

$M(n),S(n)$ ($S$는 등급환, $M$은 등급 $S$-가군): \hyperref[II.2.1.1-ko]{2.1.1}.

$\Hom_S(M,N)$ ($S$는 등급환, $M,N$은 등급 $S$-가군): \hyperref[II.2.1.2-ko]{2.1.2}.

$\mathfrak N_+,\mathfrak r_+(\mathfrak J)$
($\mathfrak J$는 $S_+$의 아이디얼): \hyperref[II.2.1.10-ko]{2.1.10}.

$S_{(f)},M_{(f)}$ ($f$는 $S_+$의 동차 원소): \hyperref[II.2.2.1-ko]{2.2.1}.

$S_{(T)},M_{(T)},S_{(\mathfrak p)},M_{(\mathfrak p)}$
($T$는 동차 원소들로 이루어진 $S_+$의 곱셈적 부분집합,
$\mathfrak p$는 $S_+$의 동차 소아이디얼): \hyperref[II.2.2.7-ko]{2.2.7}.

$\Proj(S)$ ($S$는 등급환): \hyperref[II.2.3.1-ko]{2.3.1}.

$V_+(E)$ ($E$는 $S_+$의 부분집합): \hyperref[II.2.3.2-ko]{2.3.2}.

$D_+(f)$ ($f$는 $S$의 원소): \hyperref[II.2.3.3-ko]{2.3.3}.

$\mathfrak j_+(Y)$ ($Y$는 $\Proj(S)$의 부분집합): \hyperref[II.2.3.10-ko]{2.3.10}.

$\widetilde M$ ($M$은 등급 $S$-가군): \hyperref[II.2.5.2-ko]{2.5.2}.

$\widetilde u$ ($u$는 차수 $0$인 준동형): \hyperref[II.2.5.4-ko]{2.5.4}.

$\mathcal{O}_X(n),\mathcal{F}(n)$
($X=\Proj(S)$, $\mathcal{F}$는 $\mathcal{O}_X$-가군): \hyperref[II.2.5.10-ko]{2.5.10}.

$\lambda$ (표준 준동형): \hyperref[II.2.5.11-ko]{2.5.11}.

$\mu$ (표준 준동형): \hyperref[II.2.5.12-ko]{2.5.12}.

$X^{(d)}$ ($X=\Proj(S)$, $d$는 $>0$인 정수): \hyperref[II.2.5.16-ko]{2.5.16}.

$\Gamma_*(\mathcal{F})$ ($\mathcal{F}$는 $\mathcal{O}_X$-가군,
$X=\Proj(S)$): \hyperref[II.2.6.1-ko]{2.6.1}.

$\alpha_n,\alpha,\alpha_{\mathcal{M}}$ (표준 준동형들): \hyperref[II.2.6.2-ko]{2.6.2}.

$\beta,\beta_{\mathcal{F}}$ (표준 준동형들): \hyperref[II.2.6.4-ko]{2.6.4}.

$G(\vphi),{}^a\vphi$ (표기의 남용)
($\vphi$는 등급환의 준동형): \hyperref[II.2.8.1-ko]{2.8.1}.

$\Proj(\vphi)$ ($\vphi$는 등급환의 준동형): \hyperref[II.2.8.2-ko]{2.8.2}.

$\mathcal{S}^{(d)},\mathcal{M}^{(d,k)},\mathcal{M}^{(d)},\mathcal{M}(n)$
($\mathcal{S}$는 등급 $\mathcal{O}_Y$-대수,
$\mathcal{M}$은 등급 $\mathcal{S}$-가군): \hyperref[II.3.1.1-ko]{3.1.1}.

$\Proj(\mathcal{S})$ ($\mathcal{S}$는 준연접 등급 $\mathcal{O}_Y$-대수): \hyperref[II.3.1.3-ko]{3.1.3}.

$X_f$ ($X=\Proj(\mathcal{S})$, $f$는 $Y$ 위의 $\mathcal{S}_d$의 단면): \hyperref[II.3.1.4-ko]{3.1.4}.

$\widetilde{\mathcal{M}}$ ($\mathcal{M}$은 등급 $\mathcal{S}$-가군): \hyperref[II.3.2.2-ko]{3.2.2}.

$\mathcal{O}_X(n),\mathcal{F}(n)$
($X=\Proj(\mathcal{S})$, $\mathcal{F}$는 $\mathcal{O}_X$-가군): \hyperref[II.3.2.5-ko]{3.2.5}.

$\lambda,\mu$ (표준 준동형들): \hyperref[II.3.2.6-ko]{3.2.6}.

$\Gamma_*(\mathcal{F})$ ($\mathcal{F}$는 $\mathcal{O}_X$-가군,
$X=\Proj(\mathcal{S})$): \hyperref[II.3.3.1-ko]{3.3.1}.

$\alpha_n,\alpha,\alpha_{\mathcal{M}}$ (표준 준동형들): \hyperref[II.3.3.2-ko]{3.3.2}.

\oldpage[II]{208}
\markboth{제II장}{제II장}
$\beta,\beta_{\mathcal{F}}$ (표준 준동형들): \hyperref[II.3.3.4-ko]{3.3.4}.

$G(\vphi),\Proj(\vphi)$
($\vphi$는 등급 $\mathcal{O}_Y$-대수의 준동형): \hyperref[II.3.5.1-ko]{3.5.1}.

$\Proj(u)$ ($u$는 등급 $\mathcal{O}_Y$-대수에서
등급 $\mathcal{O}_{Y'}$-대수로 가는 $g$-사상): \hyperref[II.3.5.6-ko]{3.5.6}.

$r_{\mathcal{L},\psi}$ ($\mathcal{L}$은 가역 $\mathcal{O}_X$-가군): \hyperref[II.3.7.1-ko]{3.7.1}.

$\mathbf{P}(\mathcal{E}),\mathbf{P}(E),\mathbf{P}_Y^n,\mathbf{P}_A^n$ : \hyperref[II.4.1.1-ko]{4.1.1}.

$\mathbf{P}(u)$ ($u$는 $\mathcal{O}_Y$-가군의 전사 준동형): \hyperref[II.4.1.2-ko]{4.1.2}.

$\Hyp_Y(X,\mathcal{E})$ ($\mathcal{E}$는 준연접 $\mathcal{O}_Y$-가군): \hyperref[II.4.2.5-ko]{4.2.5}.

$\varsigma$ (세그레 사상): \hyperref[II.4.3.1-ko]{4.3.1}.

$\mathcal{F}(n)$ ($\mathcal{F}$는 $\mathcal{O}_X$-가군, $X$에는
가역 $\mathcal{O}_X$-가군이 주어져 있음): \hyperref[II.4.5.1-ko]{4.5.1}.

$\varepsilon$ (항등 자기동형사상): \hyperref[II.4.5.1-ko]{4.5.1}.

$P(u,T),\sigma_i(u)$ ($u$는 유한 기저를 갖는 자유 가군의 자기준동형): \hyperref[II.6.4.1-ko]{6.4.1}.

$P(u,T),\sigma_i(u),\det u$
($u$는 정역 또는 뇌터 축소환 위의 유한형 가군의 자기준동형):
\hyperref[II.6.4.2-ko]{6.4.2} 및 \hyperref[II.6.4.7-ko]{6.4.7}.

$\sigma_i(u),\det u$ ($u$는 국소 자유 $\mathcal{A}$-가군의 자기준동형): \hyperref[II.6.4.8-ko]{6.4.8}.

$\sigma_i,\det$ (층의 준동형들): \hyperref[II.6.4.8-ko]{6.4.8}, \hyperref[II.6.4.9-ko]{6.4.9} 및 \hyperref[II.6.4.10-ko]{6.4.10}.

$\sigma_i(u),\det u$
($u$는 국소 정역 준스킴 또는 국소 뇌터 축소 준스킴 위의
유한형 $\mathcal{O}_X$-가군의 자기준동형):
\hyperref[II.6.4.9-ko]{6.4.9} 및 \hyperref[II.6.4.10-ko]{6.4.10}.

$N_{\mathcal{B}/\mathcal{A}}(f)$
($f$는 $\mathcal{A}$-대수 $\mathcal{B}$의 단면): \hyperref[II.6.5.1-ko]{6.5.1}.

$N_{\mathcal{B}/\mathcal{A}}(\mathcal{L}')$
($\mathcal{L}'$은 가역 $\mathcal{B}$-가군): \hyperref[II.6.5.2-ko]{6.5.2}.

$N_{\mathcal{B}/\mathcal{A}}(h')$
($h'$은 가역 $\mathcal{B}$-가군들의 준동형): \hyperref[II.6.5.3-ko]{6.5.3}.

$N_{X'/X}(\mathcal{L}')$
($X'$은 $X$ 위에서 유한이고, $\mathcal{L}'$은 가역 $\mathcal{O}_{X'}$-가군): \hyperref[II.6.5.5-ko]{6.5.5}.

$N_{X'/X}(h')$
($h'$은 가역 $\mathcal{O}_{X'}$-가군들의 준동형): \hyperref[II.6.5.5-ko]{6.5.5}.

$\mathcal{S}_X$ ($X=\Proj(\mathcal{S})$, $\mathcal{S}$는
$\mathcal{O}_Y$의 아이디얼들의 합): \hyperref[II.8.1.5-ko]{8.1.5}.

$S^\geq,S^\leq,S_f^\geq,S_f^\leq,
M^\geq,M^\leq,M_f^\geq,M_f^\leq$
($S$는 등급환, $M$은 등급 $S$-가군,
$f$는 $S_+$의 동차 원소): \hyperref[II.8.2.1-ko]{8.2.1}.

$\widehat S$ ($S$는 등급환): \hyperref[II.8.2.2-ko]{8.2.2}.

$\widehat M$ ($M$은 등급 $S$-가군): \hyperref[II.8.2.4-ko]{8.2.4}.

$S_{[n]},S^\natural,f^\natural$
($S$는 등급환, $f$는 $S$의 동차 원소): \hyperref[II.8.2.6-ko]{8.2.6}.

$M_{[n]},M^\natural$ ($M$은 등급 $S$-가군): \hyperref[II.8.2.8-ko]{8.2.8}.

$\widehat{\mathcal{S}}$ ($\mathcal{S}$는 등급 $\mathcal{O}_Y$-대수):
\hyperref[II.8.3.1-ko]{8.3.1}\footnote{편집자 주: 프랑스어 원문의 이 항목에는 참조 번호가 3.8.1로 인쇄되어 있으나, 연결되는 정의는 8.3.1이다.}.

$G_m$ (군 스킴): \hyperref[II.8.3.9-ko]{8.3.9}.

$\mathcal{S}_X$ ($X=\Proj(\mathcal{S})$, $\mathcal{S}$는 등급 $\mathcal{O}_Y$-대수): \hyperref[II.8.6.1-ko]{8.6.1}.

$\mathcal{S}_X,C_X,\widehat C_X,E_X,\widehat E_X$
($X=\Proj(\mathcal{S})$, $\mathcal{S}$는 $\mathcal{S}_0=\mathcal{O}_Y$인
등급 $\mathcal{O}_Y$-대수): \hyperref[II.8.6.1-ko]{8.6.1}.

$\mathcal{S}_{[n]},\mathcal{S}^{\natural}$
($\mathcal{S}$는 등급 $\mathcal{O}_Y$-대수): \hyperref[II.8.7.2-ko]{8.7.2}.

$\egaShProj_0(\mathcal{M}),\egaShProj(\mathcal{M}),\mathcal{M}_X,
\widehat{\mathcal{M}},\mathcal{M}^{\square}$
($\mathcal{M}$은 등급 $\mathcal{S}$-가군): \hyperref[II.8.12.1-ko]{8.12.1}.

$\mathcal{M}_X^\geq$
($X=\Proj(\mathcal{S})$, $\mathcal{M}$은 등급 $\mathcal{S}$-가군): \hyperref[II.8.12.5-ko]{8.12.5}.

$\mathcal{M}_{[n]},\mathcal{M}^{\natural}$
($\mathcal{M}$은 등급 $\mathcal{S}$-가군): \hyperref[II.8.12.9-ko]{8.12.9}.

$(\widetilde{\mathcal{N}})^-$
($\mathcal{N}$은 어떤 등급 $\mathcal{S}$-가군의 부분-$\mathcal{S}$-가군): \hyperref[II.8.13.1-ko]{8.13.1}.

$\overline{\mathcal{N}}$
($\mathcal{N}$은 어떤 등급 $\mathcal{S}$-가군의 부분-$\mathcal{S}$-가군): \hyperref[II.8.13.2-ko]{8.13.2}.

$\Gamma_*(\mathcal{F})$
($X=\Proj(\mathcal{S})$, $\mathcal{F}$는 등급 $\mathcal{S}_X$-가군): \hyperref[II.8.14.5-ko]{8.14.5}.

$\mathcal{S}_{X'},\Gamma'_*(\mathcal{F}'),\egaShProj'(\mathcal{M}),
\alpha',\beta'$
($X'$은 $X=\Proj(\mathcal{S})$의 열린집합, $\mathcal{F}'$은 등급
$\mathcal{S}_{X'}$-가군, $\mathcal{M}$은 등급 $\mathcal{S}$-가군): \hyperref[II.8.14.6-ko]{8.14.6}.

\endgroup

\clearpage
\oldpage[II]{209}
\markboth{제II장}{제II장}
\thispagestyle{plain}
\phantomsection
\addcontentsline{toc}{section}{용어 색인}
\section*{용어 색인}

\begingroup
\small
\setlength{\parindent}{0pt}
\setlength{\parskip}{2pt}
\newcommand{\egaTerme}[2]{\textbf{#1} : #2.\par}

\egaTerme{아핀 (사상)}{\hyperref[II.1.6.1-ko]{1.6.1}}
\egaTerme{준스킴 위에서 아핀인 준스킴}{\hyperref[II.1.2.1-ko]{1.2.1}}
\egaTerme{등급환 위의 등급대수}{\hyperref[II.2.1.2-ko]{2.1.2}}
\egaTerme{$A$-가군의 대칭 대수}{\hyperref[II.1.7.1-ko]{1.7.1}}
\egaTerme{$\mathcal{A}$-가군의 대칭 $\mathcal{A}$-대수}{\hyperref[II.1.7.4-ko]{1.7.4}}
\egaTerme{본질적으로 축소인 등급 $\mathcal{O}_Y$-대수}{\hyperref[II.3.1.12-ko]{3.1.12}}
\egaTerme{정역인 등급 $\mathcal{O}_Y$-대수}{\hyperref[II.3.1.12-ko]{3.1.12}}
\egaTerme{풍부한 가역 $\mathcal{O}_X$-가군}{\hyperref[II.4.5.3-ko]{4.5.3}}
\egaTerme{$f$에 관해 풍부한, $f$-풍부한, $Y$-풍부한 가역 $\mathcal{O}_X$-가군}{\hyperref[II.4.6.1-ko]{4.6.1}}
\egaTerme{본질적으로 정역인 등급환}{\hyperref[II.2.1.11-ko]{2.1.11}}
\egaTerme{본질적으로 축소인 등급환}{\hyperref[II.2.1.10-ko]{2.1.10}}
\egaTerme{$\mathcal{B}$-가군에 대응하는 층}{\hyperref[II.1.4.3-ko]{1.4.3}}
\egaTerme{등급 $S$-가군에 대응하는 층}{\hyperref[II.2.5.3-ko]{2.5.3}}
\egaTerme{등급 $\mathcal{S}$-가군에 대응하는 층}{\hyperref[II.3.2.2-ko]{3.2.2}}
\egaTerme{$\mathcal{O}_S$-대수에 대응하는 $S$-스킴}{\hyperref[II.1.3.1-ko]{1.3.1}}

\egaTerme{조건 (TF), 조건 (TN)}{\hyperref[II.2.7.2-ko]{2.7.2}}
\egaTerme{조건 (TF), 조건 (TN)}{\hyperref[II.3.4.2-ko]{3.4.2}}
\egaTerme{등급 $\mathcal{O}_Y$-대수로 정의되는 아핀 원뿔과 사영 원뿔}{\hyperref[II.8.3.1-ko]{8.3.1}}
\egaTerme{등급 $\mathcal{O}_Y$-대수로 정의되는 천공 아핀 원뿔과 천공 사영 원뿔}{\hyperref[II.8.3.4-ko]{8.3.4}}
\egaTerme{준스킴 $\Proj(\mathcal{S})$를 사영하는 아핀 원뿔과 사영 원뿔}{\hyperref[II.8.3.1-ko]{8.3.1}}
\egaTerme{체 위의 대수곡선}{\hyperref[II.7.4.2-ko]{7.4.2}}
\egaTerme{완비 대수곡선}{\hyperref[II.7.4.7-ko]{7.4.7}}
\egaTerme{세르 판정법}{\hyperref[II.5.2.1-ko]{5.2.1}}

\egaTerme{정역 위의 유한형 가군 자기준동형의 결정식}{\hyperref[II.6.4.2-ko]{6.4.2}}
\egaTerme{뇌터 축소환 위의 유한형 가군 자기준동형의 결정식}{\hyperref[II.6.4.7-ko]{6.4.7}}
\egaTerme{국소 자유 $\mathcal{A}$-가군 자기준동형의 결정식}{\hyperref[II.6.4.8-ko]{6.4.8}}
\egaTerme{국소 정역 준스킴 $X$ 위의 유한형 $\mathcal{O}_X$-가군 자기준동형의 결정식}{\hyperref[II.6.4.9-ko]{6.4.9}}
\egaTerme{국소 뇌터 축소 준스킴 $X$ 위의 유한형 $\mathcal{O}_X$-가군 자기준동형의 결정식}{\hyperref[II.6.4.10-ko]{6.4.10}}

\egaTerme{블로업 (준스킴)}{\hyperref[II.8.1.3-ko]{8.1.3}}
\egaTerme{정수적 (사상)}{\hyperref[II.6.1.1-ko]{6.1.1}}
\egaTerme{준스킴 위에서 정수적인 준스킴}{\hyperref[II.6.1.1-ko]{6.1.1}}
\egaTerme{정수적인 준연접 $\mathcal{O}_S$-대수}{\hyperref[II.6.1.2-ko]{6.1.2}}
\egaTerme{$\mathcal{A}$ 위에서 정수적인 $\mathcal{A}$-대수의 단면}{\hyperref[II.6.3.1-ko]{6.3.1}}

\egaTerme{사영 다발의 기본 층}{\hyperref[II.4.1.1-ko]{4.1.1}}
\egaTerme{$\mathcal{A}$-대수 안에서 환의 층 $\mathcal{A}$의 정수적 폐포}{\hyperref[II.6.3.2-ko]{6.3.2}}
\egaTerme{$\mathcal{O}_X$-대수에 대한 준스킴 $X$의 정수적 폐포}{\hyperref[II.6.3.4-ko]{6.3.4}}
\egaTerme{아핀 원뿔의 사영 폐포}{\hyperref[II.8.3.1-ko]{8.3.1}}

\oldpage[II]{210}
\markboth{A. GROTHENDIECK --- 제II장}{A. GROTHENDIECK --- 제II장}
\egaTerme{$k(y)$의 확장체 $K$에 대한 사영 다발의 $K$-유리 기하적 섬유}{\hyperref[II.4.2.6-ko]{4.2.6}}
\egaTerme{$\mathcal{O}_Y$-가군으로 정의되는 사영 다발}{\hyperref[II.4.1.1-ko]{4.1.1}}
\egaTerme{$\mathcal{O}_S$-가군으로 정의되는 벡터 다발}{\hyperref[II.1.7.8-ko]{1.7.8}}
\egaTerme{유한 (사상)}{\hyperref[II.6.1.1-ko]{6.1.1}}
\egaTerme{준스킴 위에서 유한인 준스킴}{\hyperref[II.6.1.1-ko]{6.1.1}}
\egaTerme{유한인 준연접 $\mathcal{O}_S$-대수}{\hyperref[II.6.1.2-ko]{6.1.2}}
\egaTerme{국소 자유 가군 자기준동형의 기본대칭함수}{\hyperref[II.6.4.1-ko]{6.4.1}}
\egaTerme{정역 위의 유한형 가군 자기준동형의 기본대칭함수}{\hyperref[II.6.4.2-ko]{6.4.2}}
\egaTerme{뇌터 축소환 위의 유한형 가군 자기준동형의 기본대칭함수}{\hyperref[II.6.4.7-ko]{6.4.7}}
\egaTerme{일정한 계수의 국소 자유 $\mathcal{A}$-가군 자기준동형의 기본대칭함수}{\hyperref[II.6.4.8-ko]{6.4.8}}
\egaTerme{국소 정역 준스킴 $X$ 위의 유한형 $\mathcal{O}_X$-가군 자기준동형의 기본대칭함수}{\hyperref[II.6.4.9-ko]{6.4.9}}
\egaTerme{국소 뇌터 축소 준스킴 $X$ 위의 유한형 $\mathcal{O}_X$-가군 자기준동형의 기본대칭함수}{\hyperref[II.6.4.10-ko]{6.4.10}}
\egaTerme{$\mathcal{A}$-대수의 단면의 기본대칭함수}{\hyperref[II.6.5.1-ko]{6.5.1}}

\egaTerme{등급 $\mathcal{S}$-가군의 부분-$\mathcal{S}$-가군의 동차화}{\hyperref[II.8.13.2-ko]{8.13.2}}
\egaTerme{등급환의 준동형}{\hyperref[II.2.1.2-ko]{2.1.2}}
\egaTerme{등급가군의 차수 $k$인 준동형}{\hyperref[II.2.1.2-ko]{2.1.2}}

\egaTerme{$S_+$의 아이디얼, $S_+$의 동차 소아이디얼 ($S$는 등급환)}{\hyperref[II.2.1.10-ko]{2.1.10}}

\egaTerme{저우의 보조정리}{\hyperref[II.5.6.1-ko]{5.6.1}}
\egaTerme{사영 원뿔의 무한원 자취}{\hyperref[II.8.3.3-ko]{8.3.3}}

\egaTerme{표준 사상 $G(\mathcal{E})\to\Proj(S)$, $S=\bigoplus_{n\geq0}\Gamma(X,\mathcal{L}^{\otimes n})$}{\hyperref[II.4.5.1-ko]{4.5.1}}
\egaTerme{표준 사상 $X\to\Spec(\Gamma(X,\mathcal{O}_X))$}{\hyperref[II.5.1.1-ko]{5.1.1}}
\egaTerme{세그레 사상}{\hyperref[II.4.3.1-ko]{4.3.1}}

\egaTerme{$S_+$의 멱영근기 ($S$는 등급환)}{\hyperref[II.2.1.10-ko]{2.1.10}}
\egaTerme{$\mathcal{S}_+$의 멱영근기 ($\mathcal{S}$는 준연접 등급 $\mathcal{O}_Y$-대수)}{\hyperref[II.3.1.12-ko]{3.1.12}}
\egaTerme{$\mathcal{A}$-대수의 단면의 노름}{\hyperref[II.6.5.1-ko]{6.5.1}}
\egaTerme{가역 $\mathcal{B}$-가군의 노름}{\hyperref[II.6.5.2-ko]{6.5.2}}
\egaTerme{가역 $\mathcal{B}$-가군의 단면의 노름}{\hyperref[II.6.5.3-ko]{6.5.3}}
\egaTerme{가역 $\mathcal{O}_{X'}$-가군의 노름}{\hyperref[II.6.5.5-ko]{6.5.5}}

\egaTerme{주기적 (등급대수)}{\hyperref[II.8.14.12-ko]{8.14.12}}
\egaTerme{정역 위의 유한형 가군 자기준동형의 특성다항식}{\hyperref[II.6.4.2-ko]{6.4.2}}
\egaTerme{뇌터 축소환 위의 유한형 가군 자기준동형의 특성다항식}{\hyperref[II.6.4.7-ko]{6.4.7}}
\egaTerme{아핀 (사영) 원뿔의 꼭짓점 준스킴}{\hyperref[II.8.3.3-ko]{8.3.3}}
\egaTerme{유한 표시를 갖는 등급가군}{\hyperref[II.2.1.1-ko]{2.1.1}}
\egaTerme{유한 표시를 갖는 등급 $\mathcal{S}$-가군}{\hyperref[II.3.1.1-ko]{3.1.1}}
\egaTerme{두 등급가군의 등급 텐서곱}{\hyperref[II.2.1.2-ko]{2.1.2}}
\egaTerme{사영 (사상)}{\hyperref[II.5.5.2-ko]{5.5.2}}
\egaTerme{준스킴 위에서 사영적인 준스킴}{\hyperref[II.5.5.2-ko]{5.5.2}}
\egaTerme{고유 (사상)}{\hyperref[II.5.4.1-ko]{5.4.1}}
\egaTerme{고유 (부분집합)}{\hyperref[II.5.4.10-ko]{5.4.10}}
\egaTerme{준스킴 위에서 고유한 준스킴}{\hyperref[II.5.4.1-ko]{5.4.1}}

\egaTerme{준아핀 (사상)}{\hyperref[II.5.1.1-ko]{5.1.1}}
\egaTerme{준스킴 위에서 준아핀인 준스킴}{\hyperref[II.5.1.1-ko]{5.1.1}}
\egaTerme{준아핀 (스킴)}{\hyperref[II.5.1.1-ko]{5.1.1}}

\oldpage[II]{211}
\markboth{제II장}{제II장}
\egaTerme{준유한 (사상)}{\hyperref[II.6.2.3-ko]{6.2.3}}
\egaTerme{준스킴 위에서 준유한인 준스킴}{\hyperref[II.6.2.3-ko]{6.2.3}}
\egaTerme{준사영 (사상)}{\hyperref[II.5.3.1-ko]{5.3.1}}
\egaTerme{준스킴 위에서 준사영인 준스킴}{\hyperref[II.5.3.1-ko]{5.3.1}}

\egaTerme{$S_+$의 아이디얼의 근기 ($S$는 등급환)}{\hyperref[II.2.1.10-ko]{2.1.10}}
\egaTerme{천공 사영 원뿔의 무한원 자취 위로의 표준 수축}{\hyperref[II.8.3.5-ko]{8.3.5}}

\egaTerme{$\mathcal{O}_X(-1)$의 표준 단면 ($X$는 블로업 준스킴)}{\hyperref[II.8.1.9-ko]{8.1.9}}
\egaTerme{아핀 원뿔의 영단면}{\hyperref[II.8.3.3-ko]{8.3.3}}
\egaTerme{벡터 다발의 영단면}{\hyperref[II.1.7.9-ko]{1.7.9}}
\egaTerme{아핀 (사영) 원뿔의 꼭짓점 단면}{\hyperref[II.8.3.3-ko]{8.3.3}}
\egaTerme{$\mathcal{O}_S$-대수의 스펙트럼}{\hyperref[II.1.3.1-ko]{1.3.1}}
\egaTerme{등급환의 동차 소 스펙트럼}{\hyperref[II.2.3.1-ko]{2.3.1}}
\egaTerme{준연접 등급 $\mathcal{O}_S$-대수의 동차 스펙트럼}{\hyperref[II.3.1.3-ko]{3.1.3}}

\egaTerme{(TN)-단사, (TN)-전사, (TN)-전단사 (등급가군의 준동형)}{\hyperref[II.2.7.2-ko]{2.7.2}}
\egaTerme{등급가군의 (TN)-동형사상}{\hyperref[II.2.7.2-ko]{2.7.2}}
\egaTerme{(TN)-단사, (TN)-전사, (TN)-전단사 (등급환의 준동형)}{\hyperref[II.2.9.1-ko]{2.9.1}}
\egaTerme{등급환의 (TN)-동형사상}{\hyperref[II.2.9.1-ko]{2.9.1}}
\egaTerme{(TN)-단사, (TN)-전사, (TN)-전단사 (등급 $\mathcal{S}$-가군의 준동형)}{\hyperref[II.3.4.2-ko]{3.4.2}}
\egaTerme{등급 $\mathcal{S}$-가군의 (TN)-동형사상}{\hyperref[II.3.4.2-ko]{3.4.2}}
\egaTerme{(TN)-단사, (TN)-전사, (TN)-전단사 (등급 $\mathcal{O}_Y$-대수의 준동형)}{\hyperref[II.3.6.1-ko]{3.6.1}}
\egaTerme{등급 $\mathcal{O}_Y$-대수의 (TN)-동형사상}{\hyperref[II.3.6.1-ko]{3.6.1}}
\egaTerme{$\Proj(S)$ 위의 스펙트럼 위상}{\hyperref[II.2.3.3-ko]{2.3.3}}
\egaTerme{$q$에 관해 매우 풍부한, $Y$에 관해 매우 풍부한 가역 $\mathcal{O}_X$-가군}{\hyperref[II.4.4.2-ko]{4.4.2}}
\egaTerme{유한형 (등급 $\mathcal{S}$-가군)}{\hyperref[II.3.1.1-ko]{3.1.1}}

\egaTerme{보편적으로 닫힌 (사상)}{\hyperref[II.5.4.9-ko]{5.4.9}}

\endgroup
\clearpage
\oldpage[II]{212}
\thispagestyle{empty}
\mbox{}

\clearpage
\oldpage[II]{213}
\markboth{목차}{목차}
\phantomsection
\addcontentsline{toc}{section}{원문 목차}
\section*{목차}
\thispagestyle{plain}

\begingroup
\small
\renewcommand{\arraystretch}{1.04}
\begin{longtable}{@{}p{0.09\textwidth}p{0.69\textwidth}r@{}}
\textbf{단위} & \textbf{제목} & \textbf{원문 쪽}\\
\hline
\textbf{II} & \textbf{몇몇 사상 부류에 대한 초보적 전역 연구} & \textbf{5}\\
\hyperref[section:II.1-ko]{\textsection1} & 아핀 사상 & 5\\
\hyperref[subsection:II.1.1-ko]{1.1} & $S$-준스킴과 $\mathcal{O}_S$-대수 & 5\\
\hyperref[subsection:II.1.2-ko]{1.2} & 준스킴 위에서 아핀인 준스킴 & 6\\
\hyperref[subsection:II.1.3-ko]{1.3} & $\mathcal{O}_S$-대수에 대응하는 $S$ 위에서 아핀인 준스킴 & 8\\
\hyperref[subsection:II.1.4-ko]{1.4} & $S$ 위에서 아핀인 준스킴 위의 준연접층 & 9\\
\hyperref[II.1.5-ko]{1.5} & 밑준스킴의 변경 & 12\\
\hyperref[II.1.6-ko]{1.6} & 아핀 사상 & 14\\
\hyperref[II.1.7-ko]{1.7} & 가군의 층에 대응하는 벡터 다발 & 14\\
\hyperref[section:II.2-ko]{\textsection2} & 동차 소 스펙트럼 & 19\\
\hyperref[subsection:II.2.1-ko]{2.1} & 등급환과 등급가군에 관한 일반 사항. & 19\\
\hyperref[subsection:II.2.2-ko]{2.2} & 등급환의 분수환. & 23\\
\hyperref[subsection:II.2.3-ko]{2.3} & 등급환의 동차 소 스펙트럼. & 25\\
\hyperref[subsection:II.2.4-ko]{2.4} & $\operatorname{Proj}(S)$의 스킴 구조 & 28\\
\hyperref[subsection:II.2.5-ko]{2.5} & 등급가군에 대응하는 층. & 30\\
\hyperref[subsection:II.2.6-ko]{2.6} & $\operatorname{Proj}(S)$ 위의 층에 대응하는 등급 $S$-가군. & 36\\
\hyperref[subsection:II.2.7-ko]{2.7} & 유한성 조건. & 38\\
\hyperref[subsection:II.2.8-ko]{2.8} & 함자적 성질. & 41\\
\hyperref[subsection:II.2.9-ko]{2.9} & $\operatorname{Proj}(S)$ 스킴의 닫힌 부분스킴. & 48\\
\hyperref[section:II.3-ko]{\textsection3} & 등급 대수의 층의 동차 스펙트럼 & 49\\
\hyperref[subsection:II.3.1-ko]{3.1} & 준연접 등급 $\mathcal{O}_Y$-대수의 동차 스펙트럼. & 49\\
\hyperref[subsection:II.3.2-ko]{3.2} & 등급 $\mathcal{S}$-가군에 대응하는 $\operatorname{Proj}(\mathcal{S})$ 위의 층. & 54\\
\hyperref[subsection:II.3.3-ko]{3.3} & $\operatorname{Proj}(\mathcal{S})$ 위의 층에 대응하는 등급 $\mathcal{S}$-가군. & 56\\
\hyperref[subsection:II.3.4-ko]{3.4} & 유한성 조건. & 59\\
\hyperref[subsection:II.3.5-ko]{3.5} & 함자적 성질. & 61\\
\hyperref[subsection:II.3.6-ko]{3.6} & 준스킴 $\operatorname{Proj}(\mathcal{S})$의 닫힌 부분준스킴. & 64\\
\hyperref[subsection:II.3.7-ko]{3.7} & 준스킴에서 동차 스펙트럼으로 가는 사상. & 65\\
\hyperref[subsection:II.3.8-ko]{3.8} & 동차 스펙트럼으로의 몰입 사상 판정법. & 69\\
\end{longtable}
\oldpage[II]{214}
\markboth{A. 그로텐디크 --- 제II장}{A. 그로텐디크 --- 제II장}
\begin{longtable}{@{}p{0.09\textwidth}p{0.69\textwidth}r@{}}
\hyperref[section:II.4-ko]{\textsection4} & 사영 다발. 풍부한 층 & 71\\
\hyperref[subsection:II.4.1-ko]{4.1} & 사영 다발의 정의. & 71\\
\hyperref[subsection:II.4.2-ko]{4.2} & 준스킴에서 사영 다발로 가는 사상. & 72\\
\hyperref[subsection:II.4.3-ko]{4.3} & 세그레 사상. & 76\\
\hyperref[subsection:II.4.4-ko]{4.4} & 사영 다발로의 몰입 사상. 매우 풍부한 층 & 78\\
\hyperref[subsection:II.4.5-ko]{4.5} & 풍부한 층 & 83\\
\hyperref[II.4.6-ko]{4.6} & 상대적으로 풍부한 층 & 89\\
\hyperref[section:II.5-ko]{\textsection5} & 준아핀 사상~; 준사영 사상. 고유 사상~; 사영 사상 & 94\\
\hyperref[subsection:II.5.1-ko]{5.1} & 준아핀 사상. & 94\\
\hyperref[subsection:II.5.2-ko]{5.2} & 세르 판정법. & 97\\
\hyperref[subsection:II.5.3-ko]{5.3} & 준사영 사상 & 99\\
\hyperref[subsection:II.5.4-ko]{5.4} & 고유 사상과 보편적으로 닫힌 사상. & 100\\
\hyperref[subsection:II.5.5-ko]{5.5} & 사영 사상. & 103\\
\hyperref[subsection:II.5.6-ko]{5.6} & 차우 보조정리. & 106\\
\hyperref[section:II.6-ko]{\textsection6} & 정수적 사상과 유한 사상 & 110\\
\hyperref[subsection:II.6.1-ko]{6.1} & 다른 준스킴 위에서 정수적인 준스킴. & 110\\
\hyperref[subsection:II.6.2-ko]{6.2} & 준유한 사상. & 114\\
\hyperref[subsection:II.6.3-ko]{6.3} & 준스킴의 정수적 폐포. & 116\\
\hyperref[subsection:II.6.4-ko]{6.4} & $\mathcal{O}_X$-가군의 자기준동형의 행렬식. & 120\\
\hyperref[subsection:II.6.5-ko]{6.5} & 가역층의 노름. & 125\\
\hyperref[subsection:II.6.6-ko]{6.6} & 응용: 풍부성 판정법. & 130\\
\hyperref[subsection:II.6.7-ko]{6.7} & 슈발레의 정리. & 135\\
\hyperref[section:II.7-ko]{\textsection7} & 값매김 판정법 & 138\\
\hyperref[subsection:II.7.1-ko]{7.1} & 값매김환에 대한 복습. & 138\\
\hyperref[subsection:II.7.2-ko]{7.2} & 분리성의 값매김 판정법 & 141\\
\hyperref[subsection:II.7.3-ko]{7.3} & 고유성의 값매김 판정법 & 143\\
\hyperref[subsection:II.7.4-ko]{7.4} & 대수곡선과 차원 1인 함수체 & 148\\
\hyperref[section:II.8-ko]{\textsection8} & 블로업 스킴; 사영 원뿔; 사영 폐포 & 152\\
\hyperref[subsection:II.8.1-ko]{8.1} & 블로업 준스킴 & 152\\
\hyperref[subsection:II.8.2-ko]{8.2} & 등급환에서의 국소화에 관한 예비 결과 & 157\\
\hyperref[subsection:II.8.3-ko]{8.3} & 사영 원뿔 & 162\\
\hyperref[subsection:II.8.4-ko]{8.4} & 벡터 다발의 사영 폐포 & 168\\
\hyperref[subsection:II.8.5-ko]{8.5} & 함자적 성질 & 169\\
\hyperref[subsection:II.8.6-ko]{8.6} & 꼭지점을 제거한 원뿔에 대한 표준 동형사상 & 171\\
\end{longtable}
\oldpage[II]{215}
\markboth{목차}{목차}
\begin{longtable}{@{}p{0.09\textwidth}p{0.69\textwidth}r@{}}
\hyperref[subsection:II.8.7-ko]{8.7} & 사영 원뿔의 블로업 & 173\\
\hyperref[subsection:II.8.8-ko]{8.8} & 풍부한 층과 수축 & 177\\
\hyperref[subsection:II.8.9-ko]{8.9} & 그라우어트의 풍부성 판정법: 진술 & 182\\
\hyperref[subsection:II.8.10-ko]{8.10} & 그라우에르트의 풍부성 판정법: 증명 & 184\\
\hyperref[subsection:II.8.11-ko]{8.11} & 수축의 유일성 & 189\\
\hyperref[subsection:II.8.12-ko]{8.12} & 사영하는 원뿔 위의 준연접층 & 191\\
\hyperref[subsection:II.8.13-ko]{8.13} & 부분층과 닫힌 부분스킴의 사영 폐포 & 195\\
\hyperref[subsection:II.8.14-ko]{8.14} & 등급 $\mathcal{S}$-가군에 수반되는 층에 대한 보충 & 197\\
\multicolumn{2}{@{}l}{참고문헌(계속)} & 205\\
\multicolumn{2}{@{}l}{기호 색인} & 207\\
\multicolumn{2}{@{}l}{용어 색인} & 209\\
\multicolumn{2}{@{}l}{정오표 및 추록 (목록 1)} & 217\\
\end{longtable}

\begin{flushright}
\emph{1960년 7월 22일 접수.}
\end{flushright}
\endgroup
